\subsection{Подсистема регистрации сигнала}

Подсистема регистрации выполняет оцифровку аналогового выхода синхронного
детектора и передачу данных на управляющий ПК. В установке реализованы два
альтернативных тракта регистрации: на базе промышленного модуля АЦП E-502
(основной лабораторный вариант) и на базе микроконтроллера STM32 (компактный
вариант для полевых условий).

\subsubsection{Модуль АЦП E-502}

Основным средством регистрации является модуль L-Card E-502 ---
многоканальный 16-битный АЦП промышленного класса. Модуль предоставляет
16 дифференциальных входных каналов (или 32 с общей землёй) с шестью
переключаемыми входными диапазонами от $\pm$0{,}2 до $\pm$10\,В, а также
входы внешнего тактового сигнала (DI\_SYN1) и стробирования (START\_IN);
данные передаются на ПК по интерфейсу USB\,2.0.

\paragraph{Схема подключения.}
Выход синхронного детектора подаётся на два дифференциальных входа модуля
(рис.~\ref{fig:e502-wiring}): синфазная компонента $I$ --- на контакты X3~(+)
и Y3~($-$), логический канал~2, а квадратурная компонента $Q$ --- на контакты
X4~(+) и Y4~($-$), логический канал~3. Синхронизация обеспечивается двумя
внешними сигналами: тактовый сигнал 2\,МГц поступает на вход DI\_SYN1 и
задаёт моменты отсчётов АЦП по нарастающему фронту, а стробирующий сигнал
20\,кГц --- на вход START\_IN, определяя границы окон захвата.

\begin{figure}[H]
    \centering
    \fbox{\parbox[c][0.18\textheight][c]{0.80\textwidth}{\centering
    Место для схемы подключения E-502:\\[4pt]
    Signal 1 $(\pm)$ $\to$ X3/Y3;\quad
    Signal 2 $(\pm)$ $\to$ X4/Y4;\\
    2\,МГц $\to$ DI\_SYN1;\quad
    20\,кГц $\to$ START\_IN;\quad
    GND $\to$ GND;\\
    USB 2.0 $\to$ ПК.}}
    \caption{Схема подключения модуля E-502 к синхронному детектору}
    \label{fig:e502-wiring}
\end{figure}

\paragraph{Стробированный режим захвата.}
При частоте стробирования 20\,кГц и тактовой частоте 2\,МГц каждое окно
захвата содержит $2\,000\,000 / 20\,000 = 100$ отсчётов. Для повышения
отношения сигнал/шум используется усреднение по $M = 1000$ окнам. Такой
подход эквивалентен синхронному накоплению и позволяет значительно подавить
некоррелированный шум АЦП:
\begin{equation}
    \overline{x}[n] = \frac{1}{M}\sum_{m=1}^{M} x^{(m)}[n],
    \qquad n = 0, 1, \ldots, 99,
    \label{eq:window-avg}
\end{equation}
где $x^{(m)}[n]$ --- $n$-й отсчёт в $m$-м окне. Результат --- одна
усреднённая трасса из 100 точек, соответствующая одному частотному шагу.

Два дифференциальных канала обеспечивают одновременную регистрацию $I_k$ и
$Q_k$, что непосредственно даёт комплексный S-параметр:
\begin{equation}
    S_{21}(f_k) \propto I_k + i\,Q_k.
    \label{eq:s21-iq}
\end{equation}

\paragraph{Дополнительные режимы шумоподавления.}
Помимо стробированного захвата, реализованы дополнительные режимы
шумоподавления. Первый из них использует цифровой выход DO1: переключение
между двумя состояниями позволяет чередовать измерительный и референсный захват
(с выключенным источником), а разность усреднённых трасс подавляет
систематические наводки --- по сути, это синхронное детектирование второго
порядка. Второй режим --- парное вычитание с усреднением --- чередует пары
<<сигнал/референс>>, что дополнительно снижает влияние медленных дрейфов.

\subsubsection{Микроконтроллерный АЦП на базе STM32}

Альтернативный тракт регистрации основан на микроконтроллере STM32 с
встроенным АЦП. Этот вариант обеспечивает компактность и возможность
автономной работы, однако уступает E-502 по разрядности и числу каналов.

Данные передаются на ПК по интерфейсу UART/USB со скоростью 115200\,бод.
Поддерживаются несколько форматов: текстовый ASCII (<<номер шага, значение>>),
совместимый с терминальными программами; бинарный 8-байтный, содержащий маркер
синхронизации, номер шага и два 16-битных знаковых значения ($I$ и $Q$);
бинарный формат с данными логарифмического детектора AD8317; а также
комплексный ASCII (<<номер шага, Re, Im>>) для непосредственной передачи
комплексных значений.

\paragraph{Логарифмический детектор AD8317.}
Для амплитудных измерений используется микросхема AD8317 --- широкополосный
логарифмический детектор, выходное напряжение которого пропорционально
логарифму входной СВЧ-мощности. При работе в этом режиме программная
обработка восстанавливает линейную амплитуду путём экспоненцирования
оцифрованного напряжения.

\paragraph{Сравнение трактов регистрации.}
В таблице~\ref{tab:adc-compare} приведены основные характеристики двух
трактов.

\begin{table}[H]
\centering
\caption{Сравнение трактов регистрации сигнала}
\label{tab:adc-compare}
\begin{tabular}{l c c}
\toprule
Характеристика & E-502 & STM32 \\
\midrule
Разрядность АЦП & 16\,бит & 12\,бит \\
Число каналов & 2 дифф. & 2 одн. / 1 лог. \\
Тактовый сигнал & внешний 2\,МГц & внутренний \\
Стробирование & аппаратное & программное \\
Интерфейс & USB\,2.0 & UART / USB \\
Усреднение & аппаратные окна & программное скользящее \\
Применение & лаборатория & полевые условия \\
\bottomrule
\end{tabular}
\end{table}
